\documentclass{article}
\begin{document}
\section*{高等数学 2022级期末试卷}

一、选择题（每题4分，共20分）

1. 设 $f(x)=\frac{\sin x}{x}$，则 $\lim_{x\to 0}f(x)=$ \score{4}
   A. 0 \quad B. 1 \quad C. $\infty$ \quad D. 不存在

2. 曲线 $y=x^3-3x$ 的拐点坐标为 \score{4}

3. $\int_0^1 x^2\,dx =$ \score{4}

4. 幂级数 $\sum_{n=1}^{\infty}\frac{x^n}{n}$ 的收敛半径为 \score{4}

5. 设 $A$ 为 $3\times 3$ 矩阵且 $\det(A)=2$，则 $\det(2A)=$ \score{4}

二、填空题（每题4分，共20分）

6. 函数 $f(x)=x^2$ 在 $[0,2]$ 上的平均值为 \score{4}

7. 微分方程 $y'+y=0$ 的通解为 \score{4}

8. 设 $z=x^2y+xy^2$，则 $\frac{\partial z}{\partial x}=$ \score{4}

9. 向量组 $\alpha_1=(1,0,1)^T,\alpha_2=(0,1,1)^T,\alpha_3=(1,1,2)^T$ 的秩为 \score{4}

三、计算题（每题12分，共36分）

10. 计算二重积分 $\iint_D (x+y)\,dxdy$，其中 $D$ 由 $y=x^2$ 与 $y=1$ 围成。\score{12}

11. 求矩阵 $A=\begin{pmatrix}2&1&0\0&2&1\0&0&2\end{pmatrix}$ 的特征值与特征向量。\score{12}

四、证明题（每题12分，共24分）

12. 设 $f(x)$ 在 $[a,b]$ 上连续，在 $(a,b)$ 内可导，且 $f(a)=f(b)$。证明：在 $(a,b)$ 内存在至少一个零点 $f'(\xi)=0$。\score{12}

\end{document}
